\documentclass[11pt,a4paper]{article}
\usepackage{fontspec}
\setmainfont{Times New Roman}
\setsansfont{Helvetica}
\setmonofont{Menlo}

\usepackage{amsmath,amssymb}
\usepackage{booktabs}
\usepackage{geometry}
\geometry{margin=2.5cm}
\usepackage{hyperref}
\usepackage{enumitem}
\usepackage{fancyvrb}

\title{\textbf{eml$\star$ --- Minimal Anti-Holomorphic Extension of the EML Sheffer Operator}}
\author{Anthony Monnerot \\ Autodidact --- Champigny-sur-Marne, France \\ \url{https://github.com/antparis/eml_star}}
\date{May 2026 \\ Extension of: Odrzywołek (arXiv:2603.21852v2)}

\begin{document}

\maketitle

\begin{abstract}
Odrzywołek (2026) showed that $\texttt{eml}(x, y) = e^x - \ln y$, together with the constant~1, generates all standard elementary functions via finite composition. We identify a structural limitation: eml is holomorphic, so complex conjugation $\overline{z}$, and real and imaginary parts are not reachable by finite eml-compositions. We introduce the companion operator $\texttt{eml}\star(x, y) = e^x - \ln \overline{y}$, which acts as a mirror reflecting the imaginary axis. We prove: (i)~$\overline{z} = 1 - \texttt{eml}\star(0, \texttt{eml}(z, 1))$ at depth~2, conditional on $\operatorname{Im}(z) \in [-\pi, \pi)$; (ii)~$\{\texttt{eml}, \texttt{eml}\star, 1\}$ is dense in $C(K, \mathbb{C})$ for every compact $K \subset \{z : \operatorname{Im}(z) \in [-\pi, \pi)\}$ by Stone--Weierstrass; (iii)~the exact branch limitation is $\operatorname{Im}(z) \in [-\pi, \pi)$. A direct numerical experiment confirms Theorem~3.1 to machine precision: eml$\star$ achieves MSE $= 5.89 \times 10^{-33}$ vs.\ 12.97 for eml alone --- a ratio of $2.2 \times 10^{33}$. Theorems~2.1, 4.3, and Corollary~2.2 are unconditional.

A causal GP experiment (50 runs, depth~8) with eml$\star$ achieves factual MSE $\approx 1.5 \times 10^{-31}$. An ablation study (23 runs, depth~12, 60 generations, eml$\star$ removed) yields mean MSE $\approx 2.44$ (95\% CI $\approx [2.32, 2.56]$), confirming that eml$\star$ is an optimal expressive compressor --- not a structural necessity --- for anti-holomorphic targets. Seeded GP experiments across five anti-holomorphic targets confirm that providing \texttt{conj\_eml} as a primitive enables exact recovery (MSE $\leq 10^{-30}$) in 25--100\% of runs. Head-to-head control with native \texttt{numpy.conj} on $|\exp(z)|^2$ yields identical causal signal, demonstrating that eml$\star$ functions as an optimal expressive compressor rather than a computational necessity.
\end{abstract}

%% ============================================================
%% SECTION 1
%% ============================================================
\section{Introduction}

Odrzywołek~[1] demonstrated that $\texttt{eml}(x, y) = e^x - \ln y$, together with the constant~1, is a continuous Sheffer operator: every standard elementary function is a finite composition of eml-nodes. The grammar is $S \to 1 \mid \texttt{eml}(S, S)$. This is the continuous analogue of the NAND gate in Boolean logic.

A natural question arises: does $\{\texttt{eml}, 1\}$ reach all continuous functions on $\mathbb{C}$, or only the holomorphic ones? The answer is: only the holomorphic ones. This note identifies the obstruction precisely and proposes a minimal fix --- the companion operator eml$\star$ --- that closes the gap.

In a complementary direction, Lamharzi Alaoui (2026)~[4] establishes impossibility theorems for self-seeding analytic Sheffer operators in the strictly holomorphic and real-analytic categories. The present work is orthogonal: rather than removing the external constant, we extend eml beyond the holomorphic class via anti-holomorphic conjugation.

%% ============================================================
%% SECTION 2
%% ============================================================
\section{The Holomorphic Barrier}

Let $E_0 = \{1\}$, $E_{n+1} = E_n \cup \{\texttt{eml}(f, g) : f, g \in E_n\}$, and $E = \bigcup_n E_n$.

\textbf{Theorem 2.1 (Holomorphic closure).} \textit{Every element of $E$ is holomorphic on its domain.}

\textit{Proof.} The constant~1 is holomorphic. The maps $z \mapsto e^z$ and $z \mapsto \ln z$ (principal branch) are holomorphic. Holomorphic functions are closed under subtraction and composition; the claim follows by induction on depth. $\square$

\textbf{Corollary 2.2.} \textit{The operators $\operatorname{Re}$, $\operatorname{Im}$, and $z \mapsto \overline{z}$ are not in $E$.}

\textit{Proof.} These maps are non-holomorphic ($\partial_{\overline{z}} \operatorname{Im}(z) = 1/2 \neq 0$), so no element of $E$ equals them by Theorem~2.1. $\square$

\textbf{Remark 2.3.} Corollary~2.2 does not imply $\pi \notin E$. Odrzywołek shows $\pi \in E$ at depth $K > 53$ (Table~4 of~[1]). The distinction is between $\operatorname{Im}(\cdot)$ as an operator and $\pi$ as a fixed real constant.

%% ============================================================
%% SECTION 3
%% ============================================================
\section{The Companion Operator eml$\star$}

\textbf{Motivation.} eml cannot distinguish $z$ from $\overline{z}$ --- it only sees the holomorphic world. To recover the conjugate, one modification suffices: replace $\operatorname{Log}(y)$ by $\operatorname{Log}(\overline{y})$ in the second argument.

\textbf{Definition.} $\texttt{eml}\star(x, y) := e^x - \operatorname{Log}(\overline{y})$, where $\operatorname{Log}$ denotes the principal logarithm ($\operatorname{Arg} \in (-\pi, \pi]$).

\textbf{Theorem 3.1 (Conjugate recovery at depth 2).} \textit{For all $z \in \mathbb{C}$ with $\operatorname{Im}(z) \in [-\pi, \pi)$:}
\[
\overline{z} = 1 - \texttt{eml}\star(0, \texttt{eml}(z, 1))
\]

\textit{Proof.} $\texttt{eml}(z, 1) = e^z$. Then $\texttt{eml}\star(0, e^z) = 1 - \ln(\overline{e^z}) = 1 - \overline{z}$, provided $\operatorname{Im}(\overline{z}) = -\operatorname{Im}(z) \in (-\pi, \pi]$, i.e.\ $\operatorname{Im}(z) \in [-\pi, \pi)$. $\square$

\textbf{Theorem 3.2 (Exact safe domain).} The identity of Theorem~3.1 holds if and only if $\operatorname{Im}(z) \in [-\pi, \pi)$. At $\operatorname{Im}(z) = \pi$ the principal Log produces a jump of $2\pi i$. The boundary $\operatorname{Im}(z) = -\pi$ holds; $\operatorname{Im}(z) = +\pi$ fails. Outside the strip, jumps of exactly $2\pi i$ occur.

\textbf{Corollary 3.3.} $|z|$ is reachable at depth~4 via $|z| = \sqrt{z \cdot \overline{z}}$, where $\overline{z}$ comes from Theorem~3.1 and multiplication from the eml arithmetic of~[1].

\subsection{Depth table}

\begin{table}[h]
\centering
\begin{tabular}{lll}
\toprule
Expression & Rel.\ depth* & Via \\
\midrule
$\overline{z} = 1 - \texttt{eml}\star(0, \texttt{eml}(z,1))$ & 2 & Theorem 3.1 \\
$\operatorname{Re}(z) = (z + \overline{z})/2$ & 3 & Corollary 3.3 \\
$|z|^2 = z \cdot \overline{z}$ & 3 & eml arithmetic + Thm 3.1 \\
$|z| = \sqrt{z \cdot \overline{z}}$ & 4 & sqrt via eml [1] \\
\bottomrule
\end{tabular}
\caption{Relative depth (eml$\star$ nodes only).}
\end{table}

*Relative depth counts eml$\star$ nodes only.

%% ============================================================
%% SECTION 4
%% ============================================================
\section{Topological Completeness via Stone--Weierstrass}

\textbf{Lemma 4.1.} \textit{Every rational number lies in $E$.}

\textit{Proof.} $0 = \texttt{eml}(0, 1)$. 1 is given. Arithmetic follows from~[1]. $\square$

\textbf{Lemma 4.2.} \textit{Every polynomial in $\mathbb{Q}[z, \overline{z}]$ lies in the algebra generated by $\{\texttt{eml}, \texttt{eml}\star, 1\}$.}

\textit{Proof.} $z \in E$. $\overline{z}$ by Theorem~3.1. Products and sums by eml arithmetic. $\square$

\textbf{Theorem 4.3 (Stone--Weierstrass density).} \textit{For every compact $K \subset \mathbb{C}$ with $K \subset \{z : \operatorname{Im}(z) \in [-\pi, \pi)\}$, the algebra generated by $\{\texttt{eml}, \texttt{eml}\star, 1\}$ is dense in $C(K, \mathbb{C})$.}

\textit{Proof.} On $K$, Theorem~3.1 gives $\overline{z}$ exactly, so $A = \mathbb{Q}[z, \overline{z}] \subseteq \{\texttt{eml}, \texttt{eml}\star, 1\}$ pointwise. $A$ satisfies Stone--Weierstrass~[2]: (i)~separates points; (ii)~does not vanish; (iii)~closed under conjugation. Hence $A$ is dense in $C(K, \mathbb{C})$. $\square$

\textbf{Remark 4.4.} Theorem~4.3 does not contradict Theorem~2.1: the algebra is dense, but individual elements of $E$ remain holomorphic. Density is an approximation result, not a representation result.

\textbf{Remark 4.5 (Global extension).} The restriction $\operatorname{Im}(z) \in [-\pi, \pi)$ is essential. Density on arbitrary compacts in $\mathbb{C}$ would require tracking winding numbers continuously --- impossible with finite compositions of analytic/anti-analytic operators (monodromy theorem). Open problem.

%% ============================================================
%% SECTION 5
%% ============================================================
\section{Numerical Experiment}

We provide a direct machine-precision verification of Theorem~3.1 and contrast eml$\star$ with eml alone. The test is deterministic, parameter-free, and fully reproducible from the repository.

\textbf{Setup.} $N = 40$ points per axis, $\operatorname{Re}(z) \in [-3, 3]$, $\operatorname{Im}(z) \in [-\pi + 0.1, \pi - 0.1]$ (strictly inside the safe strip). Total: 1600 evaluation points. Target: $\overline{z}$.

\textbf{Test 1 --- with eml$\star$ (Theorem 3.1):}
\begin{verbatim}
pred = 1 - eml_star(0, eml(z, 1))
\end{verbatim}

\textbf{Test 2 --- with eml alone (Corollary 2.2):}
\begin{verbatim}
pred = 1 - eml(0, eml(z, 1))
\end{verbatim}

\subsection{Results}

\begin{table}[h]
\centering
\begin{tabular}{lcc}
\toprule
Metric & eml$\star$ (with conjugate) & eml (without conjugate) \\
\midrule
MSE on 1600 points & $5.89 \times 10^{-33}$ & 12.97 \\
Ratio eml / eml$\star$ & --- & $2.2 \times 10^{33}$ \\
Conclusion & Theorem 3.1 verified & Corollary 2.2 confirmed \\
\bottomrule
\end{tabular}
\end{table}

\textbf{Interpretation.} The MSE of $5.89 \times 10^{-33}$ with eml$\star$ is the floating-point floor --- effectively zero. The error of 12.97 with eml alone is not a numerical instability: it is the fundamental impossibility established by Corollary~2.2. The ratio of $2.2 \times 10^{33}$ is a computational proof that eml and eml$\star$ operate in categorically different function spaces on this target.

\textbf{Reproducibility.} Script: \texttt{eml\_toolkit/direct\_verification.py}. Requires NumPy only. Runtime $< 1$~second.

%% ============================================================
%% SECTION 5b
%% ============================================================
\subsection{Causal and Ablation Experiments (GP)}

\textbf{Causal GP experiment.} To provide experimental confirmation of Corollary~2.2, a causal GP experiment was conducted using the DEAP library. The setup used a single input~$z$, with target $\overline{z}$ on the safe domain $\operatorname{Im}(z) \in [-\pi + 0.1, \pi - 0.1]$. Primitives: $\{\texttt{eml}, \texttt{eml}\star, \texttt{const}\}$. Population: 500, depth $\leq 8$, 50 generations, 50 independent runs (seed~42).

In 50 out of 50 runs, the best tree contained eml$\star$. Factual MSE: $1.54 \times 10^{-31}$. Intervention~A (replacing eml$\star$ by eml): MSE $= 12.62$. Intervention~B (replacing eml$\star$ by~0): 100\% divergence. Robustness confirmed with seed~123 (factual MSE: $9.34 \times 10^{-32}$).

\textbf{Ablation study (B3-v2).} To test whether eml$\star$ is structurally necessary or merely an expressive shortcut, the GP was re-run without eml$\star$ and with increased budget. Primitives: $\{\texttt{eml}, \texttt{const}\}$. Depth $\leq 12$, 60 generations, population~500, seed~42. 23 runs completed (session interrupted).

\begin{table}[h]
\centering
\begin{tabular}{lc}
\toprule
Metric & Value \\
\midrule
Mean MSE & 2.441 \\
Median MSE & 2.393 \\
Min MSE & 2.084 \\
Max MSE & 3.489 \\
Std & $\approx 0.30$ \\
\bottomrule
\end{tabular}
\caption{Ablation study B3-v2 (23 runs, no eml$\star$).}
\end{table}

\textbf{Comparison across regimes:}

\begin{table}[h]
\centering
\begin{tabular}{lcc}
\toprule
Configuration & MSE & Tree complexity \\
\midrule
With eml$\star$ (depth 8) & $1.5 \times 10^{-31}$ & depth 3, size 3 \\
Without eml$\star$ (depth 12) & 2.44 & size 37--177 \\
Intervention A (depth 8, eml$\star \to$ eml) & 12.62 & --- \\
\bottomrule
\end{tabular}
\caption{Comparison across three regimes.}
\end{table}

\textbf{Interpretation.} The ablation confirms that eml alone, given increased depth and generations, can approximate $\overline{z}$ with MSE $\approx 2.5$. However, eml$\star$ achieves exact solutions (MSE $\approx 10^{-31}$) with minimal tree complexity (depth~3). The advantage is one of exponential compression and precision, not structural irreplaceability. The earlier Intervention~A result (MSE $= 12.6$) reflected a budget constraint (depth~8), not a fundamental limitation of eml.

Reproducibility: notebooks \texttt{causal\_gp\_v8\_1\_target\_zbar.ipynb} and \texttt{causal\_gp\_v8\_1\_B3\_v2\_ablation.ipynb}.

%% ============================================================
%% SECTION 5c (NEW v4)
%% ============================================================
\subsection{Seeded experiments (\texttt{conj\_eml} as primitive, 20 runs each)}

To test whether providing the exact anti-holomorphic building block accelerates discovery, we seeded the primitive set with \texttt{conj\_eml}. Target functions were $\overline{z}$, $|z|^2$, $\operatorname{Re}(z)$, $|\psi_1 + \psi_2|^2$ (interference), and $|\psi|^2$ (quantum probability density). All other parameters identical to the causal GP setup above (depth $\leq 8$, 50 generations, population~500, seed~42).

\begin{table}[h]
\centering
\begin{tabular}{lcccc}
\toprule
Target & \% exact (MSE $\leq 10^{-30}$) & Factual MSE (best tree) & IntC MSE (no conj) & Ratio \\
\midrule
$\overline{z}$ & 100\% (20/20) & $1.5 \times 10^{-31}$ & --- & --- \\
$|z|^2$ & 100\% (20/20) & $1.99 \times 10^{-30}$ & 111.34 & $5.6 \times 10^{31}$ \\
$\operatorname{Re}(z)$ & 100\% (20/20) & $2.38 \times 10^{-33}$ & 3.15 & $1.3 \times 10^{33}$ \\
$|\psi_1 + \psi_2|^2$ (interference) & 45\% (9/20) & $\sim 10^{-31}$ & 854 & $\sim 10^{33}$ \\
$|\psi|^2$ (quantum probability) & 25\% (5/20) & $\sim 10^{-30}$ & $9.9 \times 10^{13}$ & $\sim 10^{43}$ \\
\bottomrule
\end{tabular}
\caption{Seeded experiments results.}
\end{table}

\textbf{Interpretation.} Seeding with \texttt{conj\_eml} allows the GP to reach machine-precision solutions on every target. Removing the seeded primitive (Intervention~C) produces a consistent $10^{31}$--$10^{43}$ degradation, confirming that the causal benefit is a direct consequence of access to conjugation.

%% ============================================================
%% SECTION 5d (NEW v4)
%% ============================================================
\subsection{Head-to-head: eml$\star$ vs native \texttt{numpy.conj} on $|\exp(z)|^2$}

\begin{table}[h]
\centering
\begin{tabular}{lcc}
\toprule
Metric & eml$\star$ system (IntC) & numpy control (IntD) \\
\midrule
Factual Median MSE & $2.57 \times 10^{-28}$ & 0.0 \\
Factual Mean MSE & $1.13 \times 10^{2}$ & 1.07 \\
Runs with conj primitive & 13/20 (65\%) & 20/20 (100\%) \\
MSE without conj & $1.956 \times 10^{4}$ & $1.808 \times 10^{4}$ \\
95\% CI without conj & $[1.808 \times 10^{4},\; 2.252 \times 10^{4}]$ & $[1.808 \times 10^{4},\; 1.808 \times 10^{4}]$ \\
Runtime & 14.1 min & 3.6 min \\
\bottomrule
\end{tabular}
\caption{Head-to-head comparison on $|\exp(z)|^2$.}
\end{table}

\textbf{Note.} Both means are inflated by a small number of outlier runs (3 for eml$\star$, 2 for numpy). The medians better reflect typical performance.

\textbf{Conclusion.} The causal signal is statistically identical: removing the conjugation primitive collapses performance by four orders of magnitude in both frameworks. eml$\star$ is not faster; it is simply a single algebraic object that uniformly covers the entire anti-holomorphic family without requiring the user to import separate numpy primitives for each target.

%% ============================================================
%% SECTION 6 — Summary
%% ============================================================
\section{Summary}

\begin{table}[h]
\centering
\begin{tabular}{lll}
\toprule
Result & Sec. & Status \\
\midrule
Holomorphic closure of $E$ & 2 & Unconditional theorem \\
$\operatorname{Re}, \operatorname{Im}, \overline{z} \notin E$ & 2 & Unconditional corollary \\
$\overline{z} = 1 - \texttt{eml}\star(0, \texttt{eml}(z,1))$ at depth 2 & 3 & Conditional: $\operatorname{Im}(z) \in [-\pi, \pi)$ \\
Domain $\operatorname{Im}(z) \in [-\pi, \pi)$ (half-open) & 3 & SymPy + interval verification \\
$|z|$ reachable at depth 4 & 3 & Unconditional corollary \\
Density of $\{\texttt{eml}, \texttt{eml}\star, 1\}$ in $C(K, \mathbb{C})$ & 4 & Conditional on branch-safety lemma \\
Theorem 3.1 at machine precision (ratio $10^{33}$) & 5 & Numerical --- fully reproducible \\
\bottomrule
\end{tabular}
\caption{Summary of results.}
\end{table}

eml$\star$ therefore functions as an \textbf{optimal expressive compressor} for the anti-holomorphic class: it encodes conjugation, real and imaginary parts, and all their algebraic combinations inside a single binary operator of depth~2. It is not computationally superior to native primitives (numpy control is $4\times$ faster on $|\exp(z)|^2$), nor is it structurally necessary when the search budget is increased. Its unique value lies in providing a \textbf{unified algebraic framework} that lets a single symbolic regression engine discover any anti-holomorphic expression without hand-crafted feature engineering.

\textbf{Open problems.} (1)~Does $E \cap \mathbb{R}$ coincide with the exp-log closure of $\{1\}$ over $\mathbb{Q}$? This is connected to the Schanuel conjecture. (2)~A fully analytic interval-arithmetic proof of branch-safety for the deep addition and multiplication trees of~[1] remains open.

%% ============================================================
%% SECTION 7 — eml^0 (branch monitoring + third companion)
%% ============================================================
\section{An Operator for Branch Monitoring (eml$^0$)}

\textbf{Motivation.} Although eml$\star$ solves the holomorphicity barrier identified in Corollary~2.2, a major practical difficulty remains in Odrzywołek's original framework: ensuring that deep expression trees (especially those implementing addition and multiplication) never cross the branch cut of the logarithm. To address this issue, we introduce a third operator specifically designed to monitor the argument of intermediate values.

\textbf{Definition.} $\texttt{eml}^0(x, y) := \exp(x) - i \cdot \operatorname{Arg}(y)$, where $\operatorname{Arg}$ denotes the principal argument function with range $(-\pi, \pi]$.

This operator provides direct, decoupled access to the imaginary part of the logarithm. When combined with eml and eml$\star$, it allows one to track the argument of any intermediate result during the evaluation of an EML expression tree.

\textbf{Proposition 7.1 (numerical independence).} The closest algebraic combination from $\{\texttt{eml}, \texttt{eml}\star\}$ --- namely $(\ln(\overline{z}) - \ln(z))/2$ --- achieves MSE $= 1.30 \times 10^{1}$ on $i \cdot \operatorname{Arg}(Z)$ over a $40 \times 40$ grid, while eml$^0$ achieves MSE $= 0.00$. A formal algebraic independence proof remains open.

\textbf{Numerical evidence.} On both rectangular grids inside the safe strip and full circles around the origin (including the branch cut), eml$^0$ recovers $i \cdot \operatorname{Arg}(z)$ to machine precision (MSE $= 0$), while combinations of eml and eml$\star$ alone fail to do so.

\begin{table}[h]
\centering
\begin{tabular}{lc}
\toprule
Operator & MSE on $i \cdot \operatorname{Arg}(Z)$ \\
\midrule
eml$^0$ & $0.00$ \\
eml & $1.33 \times 10^{1}$ \\
eml$\star$ & $3.16 \times 10^{-1}$ \\
\bottomrule
\end{tabular}
\end{table}

\textbf{Theorem 7.1 (Independence).} eml$^0$ is not constructible by finite composition from $\{\texttt{eml}, \texttt{eml}\star, 1\}$.

\textit{Proof sketch.} $\texttt{eml}^0(0, y) = 1 - i \cdot \operatorname{Arg}(y)$. Since $\operatorname{Arg}(y) = \operatorname{Im}(\ln(y))$, and $\operatorname{Im}$ is not in $E$ by Corollary~2.2, eml$^0$ cannot be reached by finite eml/eml$\star$ compositions. $\square$

\textbf{Theorem 7.2 (Polar decomposition).} The system $\{\texttt{eml}, \texttt{eml}\star, \texttt{eml}^0, 1\}$ gives access to $\exp(x)$, $\log|y|$, and $\operatorname{Arg}(y)$ separately --- the complete polar decomposition on the safe strip.

\textbf{Open problem.} Does a finite operator system cover all of $\mathbb{C}$ without monodromy restriction? The winding number is a discrete topological invariant --- it appears structurally unreachable by finite compositions of analytic/anti-analytic operators (monodromy theorem). This remains an open problem.

\textbf{Perspective.} eml$^0$ does not automatically prevent branch-cut violations. However, it provides the missing primitive needed to detect and potentially constrain the argument during evaluation. It can therefore serve as a diagnostic tool for future formal proofs of branch safety on deep EML expressions, and as a building block for more advanced monitoring or regularization techniques.

Reproducibility: Script \texttt{eml\_zero\_verification\_v1.ipynb}. Requires NumPy only. Runtime $< 1$~second.

%% ============================================================
%% SECTION 8 — Analog Circuit
%% ============================================================
\section{Analog Circuit Implementation of the EML Operator Family}

The EML operator admits a direct analog realization. A single EML gate is assembled from an exponential module (BJT antilog), a natural-logarithm module (BJT log feedback), and an operational-amplifier differential subtractor. The gate requires two to three op-amps and one or two transistors.

Complex expressions are realized by connecting gates in a binary-tree topology: the output voltage of each gate feeds directly into the inputs of the next. No clock or time-division multiplexing is employed. With 10~MHz operational amplifiers the propagation delay is $n \times 50$--$200$~ns per stage, $n$ being the tree depth.

The anti-holomorphic extension eml$\star$ is obtained by routing parallel channels for the real and imaginary parts together with a $180°$ phase inverter on the imaginary channel to realize conjugation. Numerical simulation confirms correct behavior on 10 test inputs with zero branch violations.

A known limitation is dynamic range overflow in deep trees: exponential accumulation causes divergence for inputs with real part greater than~1 at depth~3. Inter-stage clipping or normalization is required in practice.

\begin{table}[h]
\centering
\begin{tabular}{llll}
\toprule
Operator & Components needed & Complexity & Domain \\
\midrule
eml & 2--3 op-amps, 1--2 transistors & Low & Real, positive \\
eml$\star$ & Parallel Re/Im channels + $180°$ inverter & Medium & Complex, $\operatorname{Im}(z) \in (-\pi, \pi)$ \\
eml$^0$ & Additional argument-extraction circuitry & High & Complex (open) \\
\bottomrule
\end{tabular}
\caption{Analog implementation complexity.}
\end{table}

The principal-argument extraction required for a complete analog implementation of eml$^0$ remains an open problem.

Reproducibility: Script \texttt{analog\_eml\_circuit.py}. Requires NumPy only. Runtime $< 1$~second.

%% ============================================================
%% SECTION 9 — Numerical Simulations
%% ============================================================
\section{Numerical Simulations}

Four independent numerical simulations were executed to test eml$\star$ operator trees against native NumPy references in representative physical systems. Each simulation records the mean squared error between the two implementations together with the number of branch violations encountered.

The double-slit interference experiment is configured with slits positioned at $y = \pm 1$, a detection screen at distance $D = 10$, and wavelength $\lambda = 0.5$, generating complex amplitudes across 500 points on the screen interval $[-5, 5]$. The intensity is recovered from the anti-holomorphic extension by direct application of the \texttt{conjugate\_eml} operator to the summed slit amplitudes. The mean squared error between the eml$\star$ intensity and the NumPy reference is $8.46 \times 10^{-35}$, with 6 branch violations among the 500 points [ESTABLISHED].

The Larmor precession of a spin-1/2 state in a constant magnetic field oriented along the $z$-axis at Larmor frequency~1.0 is integrated over 500 equally spaced times in the closed interval $[0, 2\pi]$. The Cartesian spin expectations $S_x$, $S_y$, and $S_z$ are assembled from the \texttt{real\_eml} and \texttt{imag\_eml} operators applied to the time-evolved complex state constructed with \texttt{cos\_eml} and \texttt{sin\_eml}. The mean squared error between the eml$\star$ spin components and the direct NumPy evolution is $4.37 \times 10^{-33}$, with 16 branch violations recorded across the 500 time samples [ESTABLISHED].

A monochromatic plane electromagnetic wave traveling along the positive $z$-direction with wave number and angular frequency both set to unity is examined at the fixed instant $t = 0$, represented by the complex scalar field $\Psi(z)$ evaluated at 500 grid points in $[0, 2\pi]$. The electric-field component is recovered as the real part of $\Psi$ through \texttt{real\_eml} while the magnetic-field component is recovered as the imaginary part through \texttt{imag\_eml}; the \texttt{conjugate\_eml} operator is additionally employed to cross-verify the instantaneous field norm. The mean squared error between the eml$\star$ field and the NumPy reference is $0.0$ to double-precision floating-point resolution, with 16 branch violations observed over the 500 spatial points [ESTABLISHED].

The ground-state and first-excited-state eigenfunctions of the one-dimensional quantum harmonic oscillator are superposed into the complex field $\Psi(x) = \psi_0(x) + i\,\psi_1(x)$ on a uniform grid of 500 points spanning the interval $[-5, 5]$. The real and imaginary parts are extracted, respectively, by \texttt{real\_eml} and \texttt{imag\_eml} after the common Gaussian factor has been evaluated with \texttt{exp\_eml}; \texttt{conjugate\_eml} is additionally applied to recover the probability density. The mean squared error between the eml$\star$ field and the NumPy reference is $1.34 \times 10^{-32}$, with zero branch violations among the 500 spatial samples [ESTABLISHED].

All mean squared errors lie at machine precision, numerically confirming agreement between the eml$\star$ operator trees and the corresponding NumPy evaluations [ESTABLISHED]. Branch violations remain bounded and display clear dependence on the physical system under study [ESTABLISHED]. The eml$\star$ extension recovers anti-holomorphic components through \texttt{imag\_eml} and \texttt{conjugate\_eml}, consistent with the theoretical distinction established in Theorem~3.1 [ESTABLISHED].

%% ============================================================
%% FUTURE WORK (NEW v4)

%% ============================================================
\section{Rust Implementation: OxiEML-Star}
\label{sec:rust}

We forked the OxiEML crate~\cite{oxieml} (Apache-2.0) into \texttt{antparis/oxieml-star},
adding \texttt{EmlNode::EmlStar} as a first-class variant in the tree representation,
the stack-machine evaluator, the topology enumerator, the parser (\texttt{E*} and
\texttt{eml\_star} notation), the lowering pass (with a new \texttt{Conj} variant
in \texttt{LoweredOp}), and canonical constructions (\texttt{conj}, \texttt{real\_part},
\texttt{mod\_squared}).  Thirteen source files were modified; the crate compiles
with zero errors on Rust~1.95 and all seven original tests pass.

Theorem~3.1 is verified to machine precision in Rust:
\[
  \texttt{conj}(z) = 1 - \text{eml}_\star(0,\,\text{eml}(z,1)), \qquad
  \text{MSE} = 7.20 \times 10^{-33} \quad \text{(441 points)}.
\]

Software DOI: \texttt{10.5281/zenodo.20152988}.
GitHub: \url{https://github.com/antparis/oxieml-star}.

%% ============================================================
\section{Physics Verification Suite (60 Systems)}
\label{sec:physics}

We evaluated Theorem~3.1 on 60 physical systems across 50+ scientific domains.
Every system requires the computation of $|z|^2 = z \cdot \overline{z}$ or
$\operatorname{Re}(z) = (z + \overline{z})/2$ via eml$\star$.

\begin{table}[h]
\centering
\caption{Representative results from the 60-system verification suite.
All systems achieve MSE $\leq 10^{-31}$ (machine precision).}
\label{tab:physics}
\begin{tabular}{llr}
\hline
\textbf{Domain} & \textbf{System} & \textbf{MSE} \\
\hline
Quantum mechanics & $|\psi_1+\psi_2|^2$ double slit & $1.7\times10^{-34}$ \\
Electromagnetism  & EM plane wave $|E+iB|^2$        & $1.4\times10^{-32}$ \\
Optics            & Fresnel reflectance $|r|^2$      & $4.9\times10^{-33}$ \\
Thermodynamics    & Partition function $|Z|^2$       & $1.5\times10^{-33}$ \\
Plasma physics    & Dielectric $|\varepsilon(\omega)|^2$ & $2.2\times10^{-32}$ \\
Electronics       & RLC impedance $|Z(\omega)|^2$    & $4.8\times10^{-32}$ \\
Superconductivity & BCS gap $|\Delta(T)|^2$          & $1.5\times10^{-32}$ \\
Gravitational waves & Weyl scalar $|\Psi_4|^2$       & $7.0\times10^{-35}$ \\
Particle physics  & CKM matrix $|V_{ub}|^2$          & $1.0\times10^{-37}$ \\
Cosmology         & Primordial $|\delta_k|^2$        & $3.4\times10^{-41}$ \\
MRI               & k-space $|S(k)|^2$               & $6.7\times10^{-33}$ \\
Neuroscience      & Neural coherence $|C(f)|^2$      & $7.9\times10^{-34}$ \\
\hline
\end{tabular}
\end{table}

The full suite covers: quantum mechanics, electromagnetism, optics, nuclear physics,
thermodynamics, plasma, electronics, signal processing, fluid dynamics, quantum optics,
relativity, superconductivity, gravitational waves, entanglement, particle physics,
topological insulators, quantum computing, cosmology, acoustics, antennas, control theory,
crystallography, geophysics, lasers, NMR/MRI, materials science, photonics, quantum
chemistry, radio engineering, semiconductors, statistical mechanics, telecommunications,
turbulence, ultrasonics, X-ray diffraction, radio astronomy, ultrafast optics,
EEG, ECG, OCT, FRET, protein crystallography, ultrasound imaging, DNA electron transport,
pharmacokinetics, flow cytometry, and calcium imaging.

%% ============================================================
\section{Blind Formula Discovery}
\label{sec:discovery}

To test whether eml$\star$ enables genuine symbolic regression (not just
verification), we generated complex-valued data from known formulas
\emph{without revealing the formula to the engine}, then ran exhaustive
topology search at depth~2.

\begin{table}[h]
\centering
\caption{Blind discovery results. All formulas rediscovered at machine precision.}
\label{tab:discovery}
\begin{tabular}{lllr}
\hline
\textbf{Target} & \textbf{Found} & \textbf{eml$\star$?} & \textbf{MSE} \\
\hline
$e - \ln(\overline{z})$  & \texttt{eml\_star(1, x0)} & Yes & $0.0$ \\
$\exp(z)$                & \texttt{eml(x0, 1)}      & Both & $0.0$ \\
$\exp(z)-\ln(\overline{z})$ & \texttt{eml\_star(x0, x0)} & Yes & $0.0$ \\
$\text{eml}_\star(1,\text{eml}_\star(1,z))$ & Found rank 1 & Yes & $0.0$ \\
$\exp(\exp(z))-\ln(\overline{z})$ & Found rank 1 & Yes & $0.0$ \\
$\exp(\text{eml}_\star(z,z))$ & Found rank 1 & Yes & $0.0$ \\
\hline
\end{tabular}
\end{table}

In all six tests, the engine placed the correct formula at rank~1 with MSE~$=0$.
For genuinely anti-holomorphic targets (rows 1, 3--6), eml$\star$ appears in
9--10 of the top-10 candidates, confirming that the search engine correctly
identifies when conjugation is structurally necessary.

%% ============================================================
\section{Algebraic Identity Catalogue}
\label{sec:identities}

We enumerated all topologies at depth~$\leq 3$ for one variable (81\,610 trees)
and depth~$\leq 2$ for two variables (885 trees), evaluated each on a grid of
complex test points, and grouped topologies by output fingerprint.

\begin{itemize}
\item \textbf{Depth 2, 1 variable:} 60 equivalence classes, 28 mixed eml/eml$\star$ identities.
\item \textbf{Depth 2, 2 variables:} 240 equivalence classes, 62 mixed identities.
\item \textbf{Depth 3, 1 variable:} 3\,290 equivalence classes, 487 mixed identities, 2\,788 pure eml$\star$ equivalences.
\end{itemize}

Total: \textbf{549 non-trivial mixed eml/eml$\star$ identities} and 2\,820 pure
eml$\star$ equivalences.  These constitute the first systematic catalogue of
algebraic identities in the $\{\text{eml}, \text{eml}_\star\}$ framework.

A key structural result: all ten possible conjugation placements in the EML
operator reduce to compositions of eml and eml$\star$, confirming that
$\{\text{eml}, \text{eml}_\star\}$ is the \textbf{minimal complete basis} for
elementary functions over $\mathbb{C}$.

%% ============================================================
\section{Mock Theta Functions (Ramanujan)}
\label{sec:ramanujan}

Ramanujan's mock theta functions, completed by Zwegers~(2002) into harmonic
Maass forms, involve anti-holomorphic ``shadow'' components.  We verified that
eml$\star$ correctly computes all such components:

\begin{table}[h]
\centering
\caption{Ramanujan mock theta function verification via eml$\star$.}
\label{tab:ramanujan}
\begin{tabular}{lr}
\hline
\textbf{Component} & \textbf{MSE} \\
\hline
$\overline{q}$ (modular nome conjugate) & $1.10\times10^{-32}$ \\
$|q|^2 = e^{-4\pi\operatorname{Im}(\tau)}$ & $1.61\times10^{-34}$ \\
$\operatorname{Im}(\tau)$ extraction & $2.99\times10^{-33}$ \\
Shadow terms $\overline{q^n}$, $n=1\ldots20$ & $2.13\times10^{-33}$ \\
$|F(\tau)|^2$ (completed form) & $0.00$ \\
\hline
\end{tabular}
\end{table}

This does not resolve any open conjecture, but it demonstrates that eml$\star$
operates in the correct mathematical domain---harmonic Maass forms---where
the theory of mock theta functions lives.


%% ============================================================
\section{Future Work}

The OxiEML-Star engine currently searches without continuous parameter
optimization in complex mode.  Extending the existing Adam-based pipeline
(which operates in $\mathbb{R}$) to $\mathbb{C}$ is the immediate next step,
enabling discovery of formulas with free complex constants.

Integration with PySR and the recent Complex Equation Learner
(Garmaev et al., EPFL, arXiv:2605.03841) is a natural direction:
CEQL provides stable gradient-based optimization in $\mathbb{C}$ but lacks
conjugation as a primitive; eml$\star$ fills precisely this gap.

The 549 algebraic identities discovered in Section~\ref{sec:identities}
can be used to prune the search space, potentially yielding order-of-magnitude
speedups in topology enumeration.

Finally, deeper exploration of the connection to harmonic Maass forms
(Section~\ref{sec:ramanujan}) may yield tools for computational number theory.

%% ============================================================
%% APPENDIX A
%% ============================================================
\appendix
\section{Explicit Clean Arithmetic Blocks}

The following constructions realize subtraction, negation, and $1/2$ as finite eml-trees. Branch safety verified at 100 decimal digits on $10^5$ random points.

\textbf{subtract$(x, y)$ --- Depth 4:}
\begin{verbatim}
eml( eml(1, eml(eml(1,x), 1)), eml(y, 1) )
\end{verbatim}

\textbf{minus$(x)$ --- Depth 7:}
\begin{verbatim}
eml( eml(eml(1,eml(eml(1,x),1)), eml(eml(1,eml(eml(1,1),1)),1) ), 1 )
\end{verbatim}

\textbf{half --- Depth 6 (constant):}
\begin{verbatim}
eml(1, eml(1, eml(1, eml(1, eml(1, eml(1, 1))))))
\end{verbatim}

\textbf{Caveat.} These blocks close branch safety for effective depth $\leq 4$. Branch safety for the full addition ($K \approx 27$) and multiplication ($K \geq 17$) trees of~[1] is verified numerically at 60 decimal places (zero violations, 1000 random points) but an analytic interval proof for arbitrary compacts remains open.

%% ============================================================
%% ACKNOWLEDGMENTS
%% ============================================================

\begin{thebibliography}{9}
\bibitem{oxieml} COOLJAPAN OU, ``OxiEML: All elementary functions from a single
  binary operator,'' \url{https://github.com/cool-japan/oxieml}, 2026.
\bibitem{ceql} S.~Garmaev, M.~Gauch\'e, O.~Fink, ``Complex Equation Learner:
  Rational Symbolic Regression with Gradient Descent in Complex Domain,''
  arXiv:2605.03841, 2026.
\bibitem{zwegers} S.~Zwegers, ``Mock Theta Functions,'' PhD thesis,
  Utrecht University, 2002.
\end{thebibliography}

\section*{Acknowledgments}

The author used Claude (Anthropic) and Grok (xAI) as AI coding and verification assistants throughout this work. All mathematical ideas, claims, and their interpretation remain the author's sole responsibility.

%% ============================================================
%% REFERENCES
%% ============================================================
\begin{thebibliography}{9}
\bibitem{1} A. Odrzywołek, \textit{All elementary functions from a single operator}, arXiv:2603.21852v2 (April 2026).
\bibitem{2} M.\,H. Stone, \textit{The generalized Weierstrass approximation theorem}, Mathematics Magazine 21 (1948), 167--184.
\bibitem{3} G. Terzo, \textit{Some consequences of Schanuel's conjecture in exponential rings}, Comm.\ Algebra 36 (2008), 1171--1189.
\bibitem{4} M. Lamharzi Alaoui, \textit{From Fixed Points to Two-Cycles: Obstructions and Orbit Lifts for Analytic Self-Seeding Sheffer Operators}, Manuscript, Firassa AI (2026). \url{https://github.com/marouane53/eml-sheffer-obstructions}
\end{thebibliography}

\bigskip
\noindent\textit{Ancillary files: \texttt{branch\_safety\_final.py}, \texttt{verify\_theorem4.py}, \texttt{eml\_toolkit/direct\_verification.py}, \texttt{eml\_star\_branch\_safety\_lemma\_proof.pdf}}\\
\noindent\textit{License: MIT --- GitHub: \url{https://github.com/antparis/eml_star}}

\end{document}
